\documentclass[11pt]{article}
\usepackage[utf8]{inputenc}
\usepackage[T1]{fontenc}
\usepackage{geometry}
\geometry{a4paper, margin=2cm}
\usepackage{amsmath}
\usepackage{tcolorbox}
\usepackage{parskip} % espaçamento entre parágrafos

% Estilos das caixas
\tcbset{
    arc=6pt,
    boxrule=0.9pt,
    leftrule=0.6pt,
    rightrule=0.6pt,
    toprule=0.6pt,
    bottomrule=0.6pt,
    left=8pt,
    right=8pt,
    top=8pt,
    bottom=8pt,
    fonttitle=\bfseries
}

\begin{document}


\begin{center}
\Large \textbf{Revisão de Aprofundamento em Matemática} \\
\large \textbf{Análise Combinatória}
\\[0.5cm]

\textbf{Professor: Jefferson}
\end{center}

% QUESTÃO 1
\begin{tcolorbox}[title=\textbf{Questão 1}, colback=blue!6, colframe=blue!70!black]
Uma cafeteria oferece 5 tipos de café e 2 tipos de adoçante. De quantas maneiras uma pessoa pode pedir um café com exatamente 1 tipo de adoçante?
\end{tcolorbox}

\begin{tcolorbox}[title=\textbf{Solução 1}, colback=green!6, colframe=green!60!black]
\textbf{Passo 1: Identificar escolhas independentes.}\\
Há duas decisões independentes: escolha do tipo de café e escolha do adoçante.

\textbf{Passo 2: Contar as opções de cada escolha.}\\
- Opções de café: 5.\\
- Opções de adoçante: 2.

\textbf{Passo 3: Aplicar o Princípio Fundamental da Contagem.}\\
Como as escolhas são independentes, multiplicamos:
\[
5 \times 2 = 10.
\]

\textbf{Interpretação :} Existem 10 pedidos diferentes possíveis (cada tipo de café combinado com cada tipo de adoçante).\\
\emph{Comentário :} Sempre que tivermos decisões independentes e quisermos saber quantas combinações possíveis existem, multiplicamos o número de opções de cada decisão.
\end{tcolorbox}


% QUESTÃO 2
\begin{tcolorbox}[title=\textbf{Questão 2}, colback=blue!6, colframe=blue!70!black]
Uma senha eletrônica possui 4 dígitos, e cada dígito pode ser de 0 a 9. Quantas senhas diferentes existem se a repetição é permitida?
\end{tcolorbox}

\begin{tcolorbox}[title=\textbf{Solução 2}, colback=green!6, colframe=green!60!black]
\textbf{Passo 1: Entender a estrutura da senha.}\\
A senha tem 4 posições (dígitos) independentemente escolhidos de 0 a 9.

\textbf{Passo 2: Contar opções por posição.}\\
Cada posição pode assumir 10 valores (0,1,2,\dots,9).

\textbf{Passo 3: Multiplicar as opções de cada posição.}\\
Como a repetição é permitida e cada posição é independente:
\[
10 \times 10 \times 10 \times 10 = 10^4 = 10000.
\]

\textbf{Interpretação :} Há 10\,000 senhas possíveis de 4 dígitos com repetição permitida.\\
\emph{Comentário :} Quando as posições são independentes e repetição é permitida, elevamos o número de opções ao número de posições.
\end{tcolorbox}


% QUESTÃO 3
\begin{tcolorbox}[title=\textbf{Questão 3}, colback=blue!6, colframe=blue!70!black]
Usando os algarismos 1, 4, 6 e 8, quantos números de dois algarismos distintos podem ser formados?
\end{tcolorbox}

\begin{tcolorbox}[title=\textbf{Solução 3}, colback=green!6, colframe=green!60!black]
\textbf{Passo 1: Escolha do primeiro dígito.}\\
Podemos escolher qualquer um dos 4 algarismos (1,4,6,8), então 4 opções.

\textbf{Passo 2: Escolha do segundo dígito (sem repetição).}\\
Depois de usar um algarismo no primeiro dígito, restam 3 algarismos disponíveis para o segundo.

\textbf{Passo 3: Multiplicar as opções de cada posição.}\\
\[
4 \times 3 = 12.
\]

\textbf{Passo 4: Listagem (opcional) para conferir:}\\
14, 16, 18, 41, 46, 48, 61, 64, 68, 81, 84, 86 — total 12.

\textbf{Interpretação :} Existem 12 números de dois algarismos distintos formados com os algarismos dados.\\
\emph{Comentário :} Aqui usamos permutação sem repetição para posições distintas (ordem importa).
\end{tcolorbox}


% QUESTÃO 4
\begin{tcolorbox}[title=\textbf{Questão 4}, colback=blue!6, colframe=blue!70!black]
De quantas maneiras podemos organizar 5 livros diferentes em uma prateleira?
\end{tcolorbox}

\begin{tcolorbox}[title=\textbf{Solução 4}, colback=green!6, colframe=green!60!black]
\textbf{Passo 1: Interpretar o problema.}\\
Organizar livros em prateleira significa atribuir uma ordem às 5 posições disponíveis; a ordem importa.

\textbf{Passo 2: Usar permutação simples.}\\
Número de maneiras = número de permutações de 5 objetos distintos:
\[
5! = 5 \times 4 \times 3 \times 2 \times 1 = 120.
\]

\textbf{Interpretação :} Há 120 formas diferentes de organizar os 5 livros.\\
\emph{Comentário :} O fatorial n! conta todas as ordens possíveis de n objetos distintos.
\end{tcolorbox}


% QUESTÃO 5
\begin{tcolorbox}[title=\textbf{Questão 5}, colback=blue!6, colframe=blue!70!black]
Em uma caixa há 7 cartões numerados, todos diferentes. De quantas formas podemos retirar 3 cartões sem se importar com a ordem?
\end{tcolorbox}

\begin{tcolorbox}[title=\textbf{Solução 5}, colback=green!6, colframe=green!60!black]
\textbf{Passo 1: Identificar se a ordem importa.}\\
Aqui a ordem não importa (só nos interessa quais 3 cartões foram retirados).

\textbf{Passo 2: Usar combinação.}\\
Número de maneiras de escolher 3 de 7:
\[
\binom{7}{3} = \frac{7!}{3! (7-3)!} = \frac{7!}{3! \, 4!} =  \frac{7 \times 6 \times 5 \times 4!}{3! \, 4!} =\frac{7 \times 6 \times 5}{3 \times 2 \times 1} = 35.
\]

\textbf{Passo 3: Interpretação.}\\
Existem 35 subconjuntos distintos de três cartões que podemos retirar.\\
\emph{Comentário :} A combinação remove contagens duplicadas causadas por diferentes ordens do mesmo conjunto.
\end{tcolorbox}


% QUESTÃO 6
\begin{tcolorbox}[title=\textbf{Questão 6}, colback=blue!6, colframe=blue!70!black]
Em um grupo de 6 atletas, de quantas maneiras é possível escolher uma dupla para competir?
\end{tcolorbox}

\begin{tcolorbox}[title=\textbf{Solução 6}, colback=green!6, colframe=green!60!black]
\textbf{Passo 1: Ordem importa?} \\
Uma dupla não distingue quem foi escolhido primeiro; a ordem não importa.

\textbf{Passo 2: Aplicar combinação de 6 tomados 2 a 2.}\\
\[
    \binom{6}{2} =  \frac{6!}{4! (6-2)!} = \frac{6!}{4! \, 2!} =  \frac{6 \times 5 \times 4!}{4! \, 2 \times 1} = \frac{6 \times 5}{2} = 15.
\]

\textbf{Interpretação :} Há 15 possíveis duplas diferentes.\\
\emph{Comentário :} Se fosse uma dupla com posições (ex.: capitão e parceiro), usaríamos permutação; aqui não.
\end{tcolorbox}


% QUESTÃO 7
\begin{tcolorbox}[title=\textbf{Questão 7}, colback=blue!6, colframe=blue!70!black]
Uma pessoa possui 4 calças e 3 camisetas. Quantos conjuntos (calça + camiseta) ela pode montar?
\end{tcolorbox}

\begin{tcolorbox}[title=\textbf{Solução 7}, colback=green!6, colframe=green!60!black]
\textbf{Passo 1: Identificar decisões independentes.}\\
Escolha de calça (4 opções) e escolha de camiseta (3 opções) são independentes.

\textbf{Passo 2: Multiplicar as opções.}\\
\[
4 \times 3 = 12.
\]

\textbf{Interpretação :} A pessoa pode montar 12 combinações diferentes de roupa.\\
\emph{Comentário :} Cada calça pode ser combinada com cada camiseta — por isso a multiplicação.
\end{tcolorbox}


% QUESTÃO 8
\begin{tcolorbox}[title=\textbf{Questão 8}, colback=blue!6, colframe=blue!70!black]
Quantos anagramas podem ser formados com as letras da palavra \textbf{SOL}?
\end{tcolorbox}

\begin{tcolorbox}[title=\textbf{Solução 8}, colback=green!6, colframe=green!60!black]
\textbf{Passo 1: Verificar se há letras repetidas.}\\
As letras S, O, L são todas distintas.

\textbf{Passo 2: Calcular permutação simples de 3 letras.}\\
\[
3! = 3 \times 2 \times 1 = 6.
\]

\textbf{Passo 3: (Opcional) Listar os anagramas:} SOL, SLO, OSL, OLS, LSO, LOS.

\textbf{Interpretação :} São 6 anagramas possíveis.\\
\emph{Comentário :} Com letras distintas, o número de anagramas é simplesmente o fatorial do número de letras.
\end{tcolorbox}


% QUESTÃO 9
\begin{tcolorbox}[title=\textbf{Questão 9}, colback=blue!6, colframe=blue!70!black]
Serão escolhidos um líder e um assistente entre 5 estudantes. De quantas formas isso pode ser feito?
\end{tcolorbox}

\begin{tcolorbox}[title=\textbf{Solução 9}, colback=green!6, colframe=green!60!black]
\textbf{Passo 1: Reconhecer que os cargos são distintos.}\\
Líder e assistente são posições diferentes — a ordem importa.

\textbf{Passo 2: Contar as opções em sequência.}\\
- Escolha do líder: 5 opções.\\
- Após escolher o líder, escolha do assistente: 4 opções restantes.

\textbf{Passo 3: Multiplicar as escolhas.}\\
\[
5 \times 4 = 20.
\]

\textbf{Interpretação :} Existem 20 formas distintas de atribuir os dois cargos.\\
\emph{Comentário :} Cargos distintos equivalem a permutações de 2 entre 5.
\end{tcolorbox}


% QUESTÃO 10
\begin{tcolorbox}[title=\textbf{Questão 10}, colback=blue!6, colframe=blue!70!black]
Um combo de cinema inclui 3 bebidas, 4 pipocas e 2 chocolates. Quantos combos completos podem ser montados?
\end{tcolorbox}

\begin{tcolorbox}[title=\textbf{Solução 10}, colback=green!6, colframe=green!60!black]
\textbf{Passo 1: Identificar escolhas independentes.}\\
Escolha de bebida (3), escolha de pipoca (4) e escolha de chocolate (2) são independentes.

\textbf{Passo 2: Multiplicar as opções.}\\
\[
3 \times 4 \times 2 = 24.
\]

\textbf{Interpretação :} São 24 combos diferentes possíveis.\\
\emph{Comentário :} Cada escolha de bebida combina com cada escolha de pipoca e cada escolha de chocolate — multiplicamos as quantidades.
\end{tcolorbox}

\end{document}

